\begin{definitionbox}{Definition --- Real Vector Space}
\begin{definition}[Real Vector Space]\label{def:vector-spaces-real-vector-space}
A \emph{real vector space} is a set \(V\) equipped with two operations:
\[
  +:V\times V\to V,
  \qquad
  \cdot:\mathbb{R}\times V\to V,
\]
called \emph{vector addition} and \emph{scalar multiplication}, such that for
all \(u,v,w\in V\) and all \(a,b\in\mathbb{R}\):
\begin{enumerate}
  \item \(u+v=v+u\);
  \item \((u+v)+w=u+(v+w)\);
  \item there exists an element \(0_V\in V\) such that \(v+0_V=v\);
  \item for each \(v\in V\), there exists an element \(-v\in V\) such that
        \(v+(-v)=0_V\);
  \item \(a\cdot(u+v)=a\cdot u+a\cdot v\);
  \item \((a+b)\cdot v=a\cdot v+b\cdot v\);
  \item \((ab)\cdot v=a\cdot(b\cdot v)\);
  \item \(1\cdot v=v\).
\end{enumerate}
\end{definition}
\end{definitionbox}

\begin{remark*}[Standard quantified statement]
\[
\begin{aligned}
\text{\((V,+,\cdot)\) is a real vector space}
\Longleftrightarrow\;&
+:V\times V\to V
\land
\cdot:\mathbb{R}\times V\to V
\\
&{}\land
\text{the eight vector-space axioms above hold.}
\end{aligned}
\]
\end{remark*}

\begin{remark*}[Predicate reading]
\[
\text{\((V,+,\cdot)\) is a real vector space}
\]
means that the set \(V\) supports addition of vectors and multiplication of
vectors by real scalars, with the usual compatibility laws.
\end{remark*}

\begin{remark*}[Interpretation]
A real vector space is the basic algebraic setting where vectors can be added
and rescaled.  The axioms guarantee that linear combinations such as
\(a u+b v\) behave coherently.
\end{remark*}

\DefinitionalRoot

\begin{definitionbox}{Definition --- Vector Space Notation}
\begin{definition}[Vector Space Notation]\label{def:vector-spaces-notation}
Let \(V\) be a real vector space.
\begin{enumerate}
  \item For \(u,v\in V\), the expression \(u+v\) denotes \emph{vector
        addition}.
  \item For \(a\in\mathbb{R}\) and \(v\in V\), the expression \(a v\) or
        \(a\cdot v\) denotes \emph{scalar multiplication}.
  \item The element \(0_V\) denotes the \emph{zero vector} of \(V\).
  \item For \(v\in V\), the element \(-v\) denotes the \emph{additive inverse}
        of \(v\).
  \item For \(u,v\in V\), the expression \(u-v\) abbreviates
        \(u+(-v)\).
\end{enumerate}
\end{definition}
\end{definitionbox}

\begin{remark*}[Standard quantified statement]
\[
\begin{aligned}
\text{standard notation in \(V\)}
\Longleftrightarrow\;&
u+v\text{ denotes vector addition}
\\
&{}\land
a v\text{ denotes scalar multiplication}
\\
&{}\land
0_V\text{ denotes the zero vector}
\\
&{}\land
-v\text{ denotes the additive inverse of \(v\)}
\\
&{}\land
u-v=u+(-v).
\end{aligned}
\]
\end{remark*}

\begin{remark*}[Interpretation]
The subscript in \(0_V\) identifies the ambient vector space.  When the
ambient space is clear, it is standard to write \(0\) instead of \(0_V\).
\end{remark*}

\DefinitionalRoot

\begin{definitionbox}{Definition --- Vector Subspace}
\begin{definition}[Vector Subspace]\label{def:vector-spaces-subspace}
Let \(V\) be a real vector space.  A subset \(W\subseteq V\) is a
\emph{vector subspace} of \(V\) if:
\begin{enumerate}
  \item \(0_V\in W\);
  \item if \(u,v\in W\), then \(u+v\in W\);
  \item if \(a\in\mathbb{R}\) and \(v\in W\), then \(a v\in W\).
\end{enumerate}
With the operations inherited from \(V\), a vector subspace is itself a real
vector space.
\end{definition}
\end{definitionbox}

\begin{remark*}[Standard quantified statement]
\[
\begin{aligned}
\text{\(W\) is a vector subspace of \(V\)}
\Longleftrightarrow\;&
W\subseteq V
\land
0_V\in W
\\
&{}\land
\forall u,v\in W{:}\;u+v\in W
\\
&{}\land
\forall a\in\mathbb{R}\;\forall v\in W{:}\;a v\in W.
\end{aligned}
\]
\end{remark*}

\begin{remark*}[Predicate reading]
\[
\text{\(W\) is a vector subspace of \(V\)}
\]
means that \(W\) is closed under the vector-space operations inherited from
\(V\).
\end{remark*}

\begin{remark*}[Interpretation]
A subspace is a smaller linear arena inside a vector space.  The closure
conditions ensure that linear combinations of vectors in \(W\) remain in
\(W\).
\end{remark*}

\DefinitionalRoot
